% Pursuing Stacks (cotes 134-2 to 134-8): archive cote 134-7, site page 12 (PDF page 13), typescript page 454.
% typescript. Draft compared: ps6.tex:544-612.
% AI-transcribed by Claude Opus 5.5 (claude-opus-5-5, adaptive thinking, effort high), Message Batches API, batch ps3 (msgbatch_01XdFmuVVBAcgc7EFxdtCqsj), 2026-11-12; not yet checked by eye.
\scanpage{134-7/12}
\mspage{454}

between the k-modules $\Gamma_k$ ,  without having to pass to completions,
satisfying compatibility wi th augmentations and diagonal maps,  and
the extra condition  (expressing that $\Gamma(f)(0) = 0$ \xout{\unsure{)}} in exponential
notation) :
\[ f(1) = 1 \]
i.e. f reduced to component \xout{of degree} $\Gamma_k^{0}(M)\simeq k$  of degree zero is just
the identity of k with  $k\simeq\Gamma_k^{0}(M')$ .  Accordingly,  we have a less awk-%
ward \inserted{k-linearization} functor than L in \hand{(}s) ,  namely ``pointed linearization'' $L_{pt}$ :
\[ \text{(s')} \qquad M\longmapsto L_{pt}(M)\overset{def}{=}\Gamma_k(M) \quad : \quad U(k)^{\bullet}\hand{$\to$} Ab_k \unsure{\hand{)}} \qquad , \]
\xout{and} which seems to me the better candidate for describing linearization.
Thus we better rephrase now the ``question'' above in terms of \hand{(}s') rather
than \hand{(}s).  One trouble however is that the comparison map \hand{(}u)  takes \xout{value}
values in $L(X_{\hand{*}})$ ,  not $L_{pt}(X_{\hand{*}})$,  therefore we may still have to use
the ``prohibitive'' $L(X_{\hand{*}})$  as an intermediary for comparing  the complexes
$\underline{Z}^{(X_{\hand{*}}(k))}$  and $L_{pt}(X_{\hand{*}})$.  It may be noted now that,  while the first chain
\margin{(v)} complex embodies  Eilenberg-Maclane homology \hand{(v)}  (in the special case
considered above\hand{)}, the second one  $L_{pt}(X_{\hand{*}})$  (in that same case) describes
the value of the total derived functor of the familiar $\Gamma_k$ functor,  on
the ``argument'' $\pi$ placed in degree n ,  and the statement that the two
\inserted[typed]{are}\xout{should be}``the same''does sound like some standard Dold-Puppe type result
which everybody is supposed to know from the craddle  - sorry !

\begin{center}
* \qquad *\\
*
\end{center}

After another break (visit, tentative nap), still another afterthought\unsure{.}
The final shape we arrived at for the ``question'' above,  when working in
the ``pointed'' category  $U(k)^{\bullet}$  of ``pointed$^{*}$ unipotent bundles over k '' ,
was whether for any semisimplicial object $X_{*}$  in $U(k)^{\bullet}$ satisfying\hand{(}t) above
i.e.  $X_o = e$ (final object of $U(k)^{\bullet}$ ),  the two canonical maps of chain
complexes (in fact,  semisimplicial abelian groups)
\[ \text{(w)} \qquad \diagram{\begin{tikzcd} \underline{Z}^{(X_{\hand{*}}(k))} \arrow[r, "\hand{}"] & L\deleted{\illegible{1}}(X_{*}) & L_{pt}(X_{*}) \arrow[l, hook', "\hand{}"'] \end{tikzcd}} \]
are quasi-isomorphisms\inserted[typed]{(for H$_i$ with $i>0$).}\unsure{\hand{$*$}} I am not too sure yet if some extra conditions
on $X_{*}$  are not required for this to be reasonable  -  I want to review
two \over{\unsure{th}}{th}at came into my mind.

As the maps \hand{(}w) are functorial for varying $X_{\hand{*}}$ ,  it would follow

\margin{12}
